\section{Introdução}

Fluxos de desenvolvimento de software combinam filas de trabalho, políticas puxadas e cadências de planejamento~\cite{poppendieck2003lean,anderson2010kanban,reinertsen2009principles}. Em contextos ágeis, práticas inspiradas em Kanban e Scrum coexistem: um quadro visual organiza o trabalho, enquanto ritos estruturam o calendário e a transparência do progresso~\cite{schwaber2002agile,schwaber2020scrum}. Sob a ótica de pesquisa operacional, essa configuração sugere um sistema estocástico com filas, políticas de atendimento e eventos de calendário---como enfatizado pela relação $L=\lambda W$~\cite{little1961proof} e pela análise de variabilidade em sistemas com WIP limitado~\cite{hopp2008factory}; contudo, permanece difícil transmitir, em sala de aula ou em estudos exploratórios, como \emph{relações entre pares}---aqui resumidas por sinergia---interagem com a variabilidade de capacidade e com gargalos operacionais.

Simuladores educacionais comerciais, como o \emph{Kanban Simulator}~\cite{kanbanSimulator}, tornaram intuitiva a variabilidade diária e a especialização. Modelagem e simulação de processos de software são discutidas extensamente na literatura~\cite{kellner1999software,madachy2008software}, mas com foco predominante em \emph{system dynamics} ou modelos de projeto; para pesquisa e ensino avançado em PO, porém, deseja-se um modelo \textbf{enunciado em equações e parâmetros}, \textbf{reprodutível} e \textbf{separável}: deve ser possível isolar o efeito da sinergia entre dois desenvolvedores em um cartão colaborativo sem, ao mesmo tempo, alterar o \emph{backlog}, a semente aleatória ou a política de WIP.

\textbf{Objetivo.} Formalizar e apresentar um motor de simulação transparente que combina (i) fluxo Kanban diário, (ii) cadência Scrum simplificada e (iii) sinergia interpessoal na colaboração na mesma etapa e nos repasses entre etapas, permitindo experimentos A/B reprodutíveis. O texto destaca, na Seção~\ref{sec:modelo}, as \textbf{diferenças principais de abordagem} entre uma leitura \emph{sem} sinergias ($s_{ij}\equiv 0$: núcleo neutro em relação a pares) e uma leitura \emph{com} sinergias (matriz $S$ acoplada a colaboração, repasse e retrabalho).

\textbf{Contribuições.}
\begin{itemize}
  \item \textbf{Modelo operacional explícito:} tempo discreto, colunas de fluxo, capacidade por membro com componente aleatória e bônus de especialista, filas FIFO nas colunas de trabalho, e avanço entre etapas com repasses e retrabalho.
  \item \textbf{Sinergia como fator controlado:} uma matriz simétrica $s_{ij}\in[-1,1]$ entra em multiplicadores com \emph{limitação} (\emph{clamping}) na colaboração em \texttt{dev} e nos repasses entre etapas.
  \item \textbf{Calendário Scrum acoplado ao motor:} planejamento, dias úteis, revisão e retrospectiva (esta última, no protótipo, apenas como marcador de fim de sprint).
  \item \textbf{Pacote experimental mínimo:} mesma semente e \emph{backlog}, variando apenas padrões de sinergia entre pares, com exportação automática de métricas e figuras.
  \item \textbf{Artefato de software livre:} aplicação web (React/Vite) sobre o motor TypeScript, internacionalização, jogo interativo com edição de responsáveis alinhada ao motor, validação de alocação, visualizações (CFD, tempo de ciclo, relatório financeiro pedagógico) e o script \texttt{npm run paper:figures} que materializa figuras do artigo a partir do repositório.
\end{itemize}

\textbf{Organização.} A Seção~\ref{sec:relacionados} posiciona trabalhos relacionados. A Seção~\ref{sec:modelo} descreve o sistema, a notação e o algoritmo diário. A Seção~\ref{sec:ritos} liga ritos ao calendário. A Seção~\ref{sec:arquitetura} resume a arquitetura do protótipo e um fluxo visual. As Seções~\ref{sec:experimentos} e~\ref{sec:resultados} apresentam experimentos e resultados ilustrativos. A Seção~\ref{sec:discussao} discute limitações e a Seção~\ref{sec:conclusao} conclui.
